% !TEX root = ../tree_notes.tex
% Appendix A. Every formula states its conditions and points to the place
% where it is proved. Add a formula here only after its home topic exists.
% Only Chapter 1 is written so far, so most forward pointers are plain text.
\chapter{Formula Reference}\label{ch:app-formulas}

Each entry lists the conditions under which it holds and, where the chapter is
written, the place where it is proved. A formula without its conditions is a
wrong-answer verdict waiting to happen.

\section{Tree Basics}

\begin{tcolorbox}[mathbase, colback=white, colframe=primaryblue, title={Counting nodes, edges, paths (Chapter 1)}, colbacktitle=primaryblue]
\textbf{Edges of a tree.} A tree with $n$ nodes has exactly $n-1$ edges. \hfill (\thmref{c01-edges-nodes})

\textbf{Unique path.} Between any two nodes of a tree there is exactly one simple path. \hfill (\thmref{c01-unique-path})

\textbf{Height recursion.} With $\height(\nil)=-1$: for a node $v$,
$\height(v)=1+\max_{c\in\children(v)}\height(c)$, so $\height(\text{leaf})=0$ and
$\height(T)=\height(\text{root})$. \hfill (\thmref{c01-height-recursion})

\textbf{Size recursion.} $\size(v)=1+\sum_{c\in\children(v)}\size(c)$, and $\size(\nil)=0$.

\textbf{Depth vs height.} $\depth(v)$ counts edges from the root down to $v$;
$\height(v)$ counts edges from $v$ down to the deepest leaf below it. For the root,
$\depth=0$ and $\height(T)=\height(\text{root})$.
\end{tcolorbox}

\begin{tcolorbox}[mathbase, colback=white, colframe=softpurple, title={Binary tree shapes (Chapter 1)}, colbacktitle=softpurple]
\textbf{Perfect binary tree.} Height $h$ (edges from root to deepest leaf) gives
$2^{h+1}-1$ nodes and $2^{h}$ leaves; level $\ell$ holds $2^{\ell}$ nodes. \hfill (\thmref{c01-perfect-count})

\textbf{Full (proper) binary tree.} Every internal node has exactly two children, and
$\#\text{leaves}=\#\text{internal}+1$. \hfill (\thmref{c01-full-binary-leaves})

\textbf{Height bounds.} A binary tree with $n$ nodes satisfies
$\floor{\log_2 n}\le \height \le n-1$: minimum for a perfect tree, maximum for a
skewed (path) tree. \hfill (\corref{c01-height-bounds})

\textbf{Array layout.} In a $1$-indexed heap array, node $i$ has children $2i$ and
$2i+1$ and parent $\floor{i/2}$. (Chapter 4, Chapter 7.)
\end{tcolorbox}

\section{Counting Trees}

\begin{tcolorbox}[mathbase, colback=white, colframe=successgreen, title={How many trees are there? (Chapter 1)}, colbacktitle=successgreen]
\textbf{Binary tree shapes.} The number of distinct \emph{shapes} of a binary tree
with $n$ nodes is the Catalan number
$C_n=\dfrac{1}{n+1}\dbinom{2n}{n}$, with $C_0=1,1,2,5,14,42,\dots$. \hfill (\thmref{c01-catalan-shapes})

\textbf{Labeled trees (Cayley).} The number of distinct trees on $n$ \emph{labeled}
vertices is $n^{\,n-2}$. \hfill (\thmref{c01-cayley})

\textbf{Labeled forests / rooted trees.} Rooting a labeled tree at any of its $n$
vertices gives $n^{\,n-1}$ labeled rooted trees.
\end{tcolorbox}

\section{Distances and Ancestors}

\begin{tcolorbox}[mathbase, colback=white, colframe=accentorange, title={Paths via LCA (Chapters 9 and 11)}, colbacktitle=accentorange]
\textbf{Distance from LCA.} For nodes $u,v$ of a rooted tree,
$\dist(u,v)=\depth(u)+\depth(v)-2\,\depth(\lca(u,v))$.

\textbf{Diameter.} $\diam(T)=\max_{u,v}\dist(u,v)$; it equals the largest
$\depth(a)+\depth(b)$ over pairs whose LCA is the branching node, computable by two
DFS/BFS sweeps or by a tree DP.

\textbf{On a path.} $u$ is an ancestor of $v$ iff $\tin(u)\le\tin(v)$ and
$\tout(v)\le\tout(u)$ in the Euler (DFS-order) tour. (Chapter 11.)
\end{tcolorbox}

\begin{notebox}[Convention used in this book]
$\height$ and $\depth$ are counted in \textbf{edges}, so a single node has height
$0$ and $\height(\nil)=-1$. Some references count in \emph{nodes} (a single node has
height $1$); the two differ by exactly $1$. Always check which convention a problem
uses before you submit.
\end{notebox}
